%! @@TITLE@@ — Unit @@UNITINT@@, Lesson @@LESSONID@@ — Guided Notes & Practice
% UNITS 3+ PACKET-FIRST shape (LESSON_SHAPE.md §0a, 2026-11-14). THE in-class document, 45
% minutes: THREE ROUNDS, then Closing Practice. The projected deck mirrors this page frame for
% frame. In Unit 3 the rounds are the lesson's three merged skills; from Unit 4 on, three
% climbing cases of the one skill. FOUR PAGES IS A CEILING, blank and key page for page.
% Each round, in order:
%   1. DEFINITION  — the term in bold, a COMPLETE printed sentence (nothing to fill in), a
%                    pre-drawn display where one helps.
%   2. EXAMPLE     — the I DO: the problem printed, a work block the student fills as the
%                    teacher works it (the key prints every step with its reason).
%   3. YOUR TURN   — the ladder \rung{1} / \rung{2} / \rung{3}: the example's shape on fresh
%                    numbers, a work block each. ★★★ is for whoever finishes early.
%   4. CHECK       — Show me / Find the mistake runs on the deck; the page carries only a
%                    "My answer" blank (Show me) or Student A's work + \writelines{2}.
%                    The LAST round's check is the crux.
% Then CLOSING PRACTICE: four mixed problems (≥1 per round) + a bonus (+1) in a student-life
% context, "score yourself out of 4".
% No vocabbox, no objectivebox, no guidednotes table, no \answerspace. NO SKETCH QUESTIONS.
% \boxguard before every notesbox. ≤3 minutes of teacher talk a round (§1 engagement rules).
\documentclass[12pt]{article}
\usepackage{@@PREFIX@@-article}
\usepackage{@@PREFIX@@-boxes}
\newcommand{\rung}[1]{{\color{goldacc}\ifcase#1\or$\star$\or$\star\star$\or$\star\star\star$\fi}}

\begin{document}

\pageheader{Unit @@UNITINT@@, Lesson @@LESSONID@@}{Guided Notes \& Practice}
% NAMESTRIP: no \namedateperiod here — the student writes their name once, on the cover.

\vspace{0.10in}

\boxguard[20]
\begin{notesbox}{Round 1 --- TODO: the first skill}
\small
\textbf{TODO term.} TODO: the definition as one complete printed sentence.

\vspace{6pt}
\textbf{Example.} TODO: the problem.
\begin{work}
  % TODO — identical in the key; the key prints the steps.
  x &= x
\end{work}

\vspace{6pt}
\textbf{Your turn.} \textit{Start at \rung{1} and climb.}

\rung{1}\; TODO
\begin{work}
  x &= x
\end{work}

\rung{2}\; TODO
\begin{work}
  x &= x
\end{work}

\rung{3}\; TODO
\begin{work}
  x &= x
\end{work}

\textbf{Show me.} My answer: \blank{1.0cm}
\end{notesbox}

% TODO: Round 2 — the same four parts.

% TODO: Round 3 — the same four parts; its check (often Find the mistake) is the crux:
%   \textbf{Find the mistake.} Student A wrote: ... Which step is wrong? Fix it. \writelines{2}

\boxguard[20]
\begin{notesbox}{Closing Practice --- score yourself out of 4}
\small
\begin{enumerate}[leftmargin=1.9em, itemsep=8pt]
  \item TODO (Round 1) \blank{3.0cm}
  \item TODO (Round 2) \blank{3.0cm}
  \item TODO (Round 3) \blank{3.0cm}
  \item TODO (mixed) \blank{3.0cm}
\end{enumerate}
\textbf{Bonus (+1).} TODO: a student-life context. \blank{3.0cm}
\end{notesbox}

\end{document}
